% !TeX root = ../../main.tex
\question{[عنوان سوال اول]}
\label{sec:question1}

[خلاصه مسئله، فرض‌ها و روش حل خود را در این قسمت بنویسید.]

\subsection{[عنوان بخش اول]}

الف) [پاسخ خود را همراه با مراحل حل و استدلال بنویسید.]

برای نوشتن اصطلاح انگلیسی در متن فارسی از دستور
\lr{\texttt{\textbackslash lr\{...\}}}
استفاده کنید؛ برای مثال: \lr{Variational Autoencoder (VAE)}.
رابطه زیر نمونه‌ای از درج فرمول و ارجاع به آن است:

\begin{equation}
  \mathcal{L}(\theta,\phi;x)
  = \mathbb{E}_{q_\phi(z\mid x)}[\log p_\theta(x\mid z)]
    - D_{\mathrm{KL}}\!\left(q_\phi(z\mid x)\,\|\,p(z)\right).
  \label{eq:elbo}
\end{equation}

در رابطه~\eqref{eq:elbo}، جمله اول به بازسازی داده و جمله دوم به
منظم‌سازی توزیع نهفته مربوط است.

\subsection{[عنوان بخش دوم]}
\subsubsection{[عنوان زیربخش]}

[پاسخ زیربخش، روابط مورد استفاده و نتیجه نهایی را در این قسمت بنویسید.]
برای نمونه، مراحل پاسخ می‌توانند به این صورت باشند:

\begin{enumerate}
  \item تعریف فرض‌ها و نمادهای مورد استفاده؛
  \item ارائه مراحل حل یا الگوریتم پیشنهادی؛
  \item تحلیل نتیجه و بیان محدودیت‌ها.
\end{enumerate}
